% Completed source-scope proofs in this authoring checkpoint.

\begin{annexitem}{annex:p17.htn-finite-derivation}{thm:p17-htn-finite-derivation}{Finite decomposition depth and the solution operator}
\begin{proof}
Induct on \(r\). At zero, the definition is exactly primitive completion.
If the statement holds at \(r\), membership at \(r+1\) either already holds
at \(r\), or holds at \(r\) for one immediate reduction \(d'\). The former
has a derivation of length at most \(r\); prefixing the latter derivation by
the selected reduction gives length at most \(r+1\).

Conversely, a derivation of length at most \(r+1\) either has length at most
\(r\), or has an initial reduction followed by at most \(r\) further steps.
The induction hypothesis and the defining union give membership in either
case. Taking the union over finite \(r\) proves the final assertion. The
definition adds no limit-stage execution for an infinite derivation.
\end{proof}
\end{annexitem}

\begin{annexitem}{annex:p17.htn-example}{thm:p17-htn-example}{A protected execution and failures of two different kinds}
\begin{proof}
Rename the method parameter to a fresh variable \(v\). Unifying its head
\(\mathit{job}(v)\) with the initial task \(\mathit{job}(w)\) gives
\(\theta=\{v\mapsto w\}\). All method literals and actions receive this
same substitution, producing the ground abbreviations in the example.
Decomposing \(j\) replaces it by the fresh labels \(a,b\) and conjoins
\(a\prec b\), \((a,g,b)\) and \((b,c)\). Label \(z\) remains present.
For the order \(a,b,z\), the trace is
\[
 \{g\},\quad\{g,r\},\quad\{g,r,c\},\quad\{r,c\}.
\]
The setup and run preconditions hold, and disturbance is always executable.
The protected interval has indices \(1\leq e<2\), so only the state
\(\{g,r\}\) is checked and it satisfies \(g\). The after-condition at
\(b\) holds because \(c\in\{g,r,c\}\). This is a primitive completion
after one permitted decomposition, hence a solution.

For order \(a,z,b\), the executable trace is
\[
 \{g\},\quad\{g,r\},\quad\{r\},\quad\{r,c\}.
\]
Run still has its required readiness atom, and the final completion atom is
true. But the protected interval now has \(1\leq e<3\); its state at
\(e=2\) lacks \(g\). The primitive completion condition fails.

For the job-only network, its sole method generates exactly the two primitive
occurrences setup and run, ordered in that order. There is no further
decomposition and no primitive-task insertion rule. Shortcut is executable
from \(I\) and gives \(\{g,c\}\), but its one-action sequence is not a
completion of the method-derived network. The finite-derivation characterization
therefore excludes it from the solution set.

If the only method for loop replaces loop by one loop occurrence, induction
on the number of reductions shows that a nonprimitive occurrence remains
after every finite derivation. Its completion set is empty, and so is the
solution set. With a second method to a dummy primitive, one decomposition
gives an executable completion. Choosing the recursive method repeatedly
still produces an infinite branch, proving the stated search limitation.
\end{proof}
\end{annexitem}

\begin{annexitem}{annex:p17.htn-semidecision}{thm:p17-htn-semidecision}{An effective enumeration of finite decompositions}
\begin{proof}
There are finitely many constants and no function symbols. Each finite set
of variables therefore has finitely many ground assignments. Enumerate the
groundings of the initial network and, at each reduction, the ground instances
of a matching method. Use a canonical sequence of fresh labels, so infinite
choices of names do not introduce infinite branching. The finite method set
and the finite number of task occurrences make each reduction node finitely
branching.

At every depth enumerate all such reductions. For each primitive network,
enumerate its finitely many total orderings and simulate the actions. Every
state and constraint check is effective; false preconditions or a false Boolean
constraint reject that candidate. Halt with its action sequence when a
completion is found. An empty primitive network is checked directly.

Every reported sequence has the required finite derivation and primitive
completion. Conversely, any solution has a finite derivation by
Theorem~\ref{thm:p17-htn-finite-derivation}. Ground values used in its final
completion give compatible choices for every variable along that derivation.
The enumeration reaches those choices, that derivation and that ordering at
a finite stage. Thus every solvable problem is eventually accepted. Without
a solution there may be arbitrarily deep compound decompositions, so this
procedure provides no general negative stopping rule.
\end{proof}
\end{annexitem}

\begin{annexitem}{annex:p17.htn-cfl-intersection}{thm:p17-htn-cfl-intersection}{Alternating checks of finite accepting computation histories}
\begin{proof}
Acceptance of an input by a deterministic Turing machine is undecidable:
a machine can simulate any given computation and enter an accepting state
exactly when that computation halts, reducing the halting problem
(Theorem~\ref{thm:halting-undecidable}). We effectively construct two
context-free languages whose intersection records the accepting histories
of a given machine on a given input.

Encode a configuration as a finite tape word with boundary markers and
exactly one marked symbol \((q,a)\), recording the control state and the
scanned tape symbol. An accepting finite computation uses only finitely
many cells. Its configurations can therefore all be padded with blanks to
one common width, with no transition crossing the boundary. Valid padded
configurations form a regular language. Initial configurations consist of
the specified input and a blank tail, and final accepting configurations
have their marker in an accepting control state; these conditions are
regular as well.

The relation between two consecutive configurations of equal width is
recognizable by a finite automaton reading their aligned pairs of symbols.
For example, a right move \((q,a)\mapsto(q',b,R)\) changes the aligned
pairs only in the adjacent window
\[
 ((q,a),b),\quad(c,(q',c)),
\]
where \(c\) is the old symbol of the next cell; everywhere else the pairs
are diagonal \((d,d)\). A left move uses the reversed-position window,
and a stay move changes only the marked cell. Take the finite union of
these patterns over the transition table and intersect with the regular
one-head and boundary conditions. This gives an effective finite automaton
for one legal step. Reversing its accepted pair words is also regular:
reverse its edges and exchange initial and final roles in a finite
nondeterministic automaton.

Consider words of the form
\[
 C_0\#C_1^R\#C_2\#C_3^R\#\cdots\#C_t^{\epsilon_t},
\]
where the separator is fresh, even configurations are written forwards,
odd configurations backwards, and \(\epsilon_t\) denotes that parity
choice. A finite-state condition checks the syntactic configuration form,
that the first is an initial configuration for the given input, and that
the final one is accepting. It does not check transitions or common width.

A pushdown automaton for the first language checks pairs
\((C_0,C_1),(C_2,C_3),\ldots\). Push the first configuration of a pair,
then pop it while reading the reversed next configuration. This aligns
their symbols in reverse order, so the reversed finite transition checker
verifies a legal step. Emptying the stack exactly at the configuration
boundary also verifies equal widths. Reset at the next pair. If there is
an unpaired final configuration, a nondeterministic branch can read it
without a pair check, but accepts only if it is the last configuration.

A second pushdown automaton skips \(C_0\) and checks pairs
\((C_1,C_2),(C_3,C_4),\ldots\). Pushing the written reversal \(C_1^R\)
and popping while reading \(C_2\) aligns the configurations in forward
order, so it uses the forward transition checker. It similarly handles
a possible unpaired final configuration. Both machines run the regular
syntax/initial/final check in their finite control. Their intersection
therefore checks every adjacent transition, forces all adjacent widths
to agree, and admits exactly finite accepting histories. A history of
one accepting initial configuration is permitted by both machines.

For completeness of the grammar construction, convert each pushdown
machine effectively to empty-stack acceptance and normalize its stack
replacement operations to replace one symbol by zero, one or two symbols.
Introduce grammar variables \([p,A,q]\), denoting strings that take the
machine from control state \(p\) with top symbol \(A\) to state \(q\)
just after removing \(A\), leaving the lower stack untouched. A transition
reading \(a\), with \(a\) possibly empty, which pops \(A\) from state
\(p\) into \(q\), gives \([p,A,q]\to a\). Replacing \(A\) by
\(B\) and entering state \(r\) gives
\([p,A,q]\to a[r,B,q]\). Replacing it by \(BC\), with \(B\) on top,
gives, for every intermediate state \(s\),
\[
 [p,A,q]\ \to\ a[r,B,s][s,C,q].
\]
The start variable selects the initial control/bottom-stack symbol and
an accepting final control state. Induction on these productions shows
that every generated word has the stated stack computation. Conversely,
decompose an accepting stack computation at the moments its newly pushed
symbols are removed; induction on its finite computation gives a grammar
derivation. Thus the two effective grammars generate exactly the two
pushdown languages just described.

Their intersection is nonempty exactly when the machine accepts its input,
so intersection nonemptiness is undecidable. All history words are nonempty.
Encode each alphabet symbol by a distinct binary block of one fixed length.
The resulting word map is injective, so equality and intersection
nonemptiness are preserved. Replace terminals in productions by their
blocks, compute nullable nonterminals and omit their occurrences in all
possible ways, retaining the nonempty language; an omitted occurrence
can be re-expanded by its empty derivation, proving preservation.
Eliminate unit productions by their finite transitive closure, introduce
one new nonterminal for each terminal occurring in a longer right-hand
side, and split longer right-hand sides into binary chains. These
effective transformations preserve the nonempty generated language and
leave precisely rules \(A\to BC\) and \(A\to a\). They give the
claimed binary, nonempty, normal-form restriction.
\end{proof}
\end{annexitem}

\begin{annexitem}{annex:p17.htn-undecidable}{thm:p17-htn-undecidable}{Two grammars and an alternating production-consumption protocol}
\begin{proof}
The source reduction uses the undecidability of nonemptiness of the
intersection of two context-free languages
\cite[Appendix, Theorem~1]{ErolHendlerNau1996HTNComplexity}; this is the
specific formal-language premise from \cite{HopcroftUllman1979HTNPremise},
proved here in Annex~\ref{annex:p17.htn-cfl-intersection}.
Use two grammars over a common binary alphabet, with disjoint nonterminal
sets, in Chomsky normal form, and with no empty word. These restrictions
preserve the problem: a fixed-length injective binary encoding of a finite
alphabet preserves equality of generated words; ordinary normal-form conversion
gives rules \(A\to BC\) or \(A\to a\). Empty-word membership is decidable.
If both grammars contain the empty word, return a trivial solvable instance;
otherwise remove empty words, which cannot contribute to their intersection.

For each terminal \(a\), introduce a state atom \(p_a\) and two primitive
actions. Action \(a_1\) has precondition \(\mathit{turn}\), adds \(p_a\)
and deletes turn. Action \(a_2\) has preconditions \(p_a,\neg\mathit{turn}\),
deletes \(p_a\) and adds turn. Initially only turn is true. For every
nonterminal introduce a compound symbol. A grammar rule \(A\to BC\)
gives an ordered two-child method for the corresponding compounds; a rule
\(A\to a\) gives the one-child method to \(a_1\) in grammar 1 and to
\(a_2\) in grammar 2.

The initial network has exactly two unordered compound occurrences for the
start symbols, labelled \(u,v\), and the constraint
\[
 (\operatorname{last}[u,v],\mathit{turn}).
\]
This tests turn after the final primitive action. It imposes no order between
the two roots. All written method bodies have at most two tasks and are
totally ordered; the initial network also has at most two tasks.

Inductively, whenever turn is true, every \(p_a\) is false. The next
nondummy action must be some \(a_1\), leaving just \(p_a\) true and turn
false. The only executable next action is then \(a_2\), which consumes
that atom and restores the invariant. Hence every executable sequence
alternates matching terminal actions, starting with grammar 1. The final
constraint forces it to finish with grammar 2.

A finite decomposition of each root is precisely a derivation in its grammar.
The ordered methods preserve the left-to-right order of its terminal yield,
while the roots' primitive descendants may interleave. A completed plan
therefore projects to the same terminal word in both grammars. Conversely,
given a common nonempty word and its two parse trees, decompose the roots
according to those trees and interleave the two terminal yields in matching
pairs. The resulting execution meets all local ordering constraints and the
final turn condition.

Thus the HTN instance has a solution exactly when the grammar intersection
is nonempty. The construction has polynomial size and meets the stated
two-task restrictions, proving undecidability. The semidecision theorem
provides the positive enumeration, so plan existence is strictly semidecidable.
\end{proof}
\end{annexitem}

\begin{annexitem}{annex:p17.htn-fixed-undecidable}{thm:p17-htn-fixed-undecidable}{Encoding grammar rules as initial-state data in a fixed domain}
\begin{proof}
Use the fixed terminal alphabet \(\{a,b\}\), atoms \(p_a,p_b,\mathit{turn}\),
the four production/consumption actions of
Annex~\ref{annex:p17.htn-undecidable}, and one effect-free dummy action.
Reserve distinct constants \(A_1,B_1,A_2,B_2,D\) and one unary compound
symbol \(t\). Four methods map \(t(x)\) to the appropriate terminal action
under the equality constraint \(x=A_i\) or \(x=B_i\). A fifth maps
\(t(x)\) to the dummy action under \(x=D\).

A static ternary predicate \(R\) records grammar rules. The sixth method
decomposes \(t(x)\) into ordered tasks \(t(y),t(z)\), requiring
\(R(x,y,z)\) after the first child. No operator changes \(R\), so this
condition is equivalent to requiring the corresponding initial fact. The
domain consists of these five operators and six methods, independently of
the grammars.

For the two normalized grammars, introduce distinct constants for all their
nonterminals, disjoint from the reserved constants. A rule \(X\to YZ\)
becomes the fact \(R(X,Y,Z)\). A terminal rule \(X\to a\) in grammar
\(i\) becomes \(R(X,A_i,D)\), and likewise for \(b\). Add turn to the
initial state, and start with the two unordered tasks \(t(S_1),t(S_2)\),
again requiring turn after their last primitive descendant.

The static facts permit exactly the grammar productions. Equality constraints
prevent nonterminal constants from being confused with terminal or dummy
constants. Each terminal yield carries the grammar's correct subscript;
the inserted dummy actions change neither the state nor the grammar word.
After erasing these dummies, the alternating argument of the preceding annex
again forces the two yielded words to agree. Conversely, two parses of a
common word yield a valid interleaving, with their dummy occurrences inserted
in the local orders required by the parses.

Only constants, initial facts and the initial task parameters depend on the
input grammars. The construction is effective and polynomial. Thus this one
fixed domain has undecidable plan existence; semidecidability still follows
from the general enumeration.
\end{proof}
\end{annexitem}

\begin{annexitem}{annex:p17.htn-levels}{thm:p17-htn-levels}{A finite search bound from a decreasing ground-task rank}
\begin{proof}
Ground the finite function-free description. Every edge from a method head
to a child strictly decreases an integer between zero and \(k\). Hence no
root-to-leaf decomposition path has more than \(k\) edges. Let \(r\) be
the initial number of task occurrences and \(B\geq1\) the maximum number
of children in a method. Every decomposition forest then has at most
\[
 r(1+B+\cdots+B^k)
\]
nodes. There are finitely many choices of ground methods at each node, so
there are finitely many such labelled forests up to the canonical renaming
of occurrence labels. Their primitive leaves also admit only finitely many
orders. Execute and check every candidate. Acceptance is sound and complete
by the finite-derivation theorem, and exhaustive failure is now a decision.

The finite directed graph of ground tasks, with an edge for every possible
method-head/child occurrence, also makes the acyclicity promise checkable.
A directed cycle contradicts any strictly decreasing level assignment.
If the graph is acyclic, assign a task the length of a longest downward path.
Primitive tasks have no outgoing method edges and receive zero. These levels
give an appropriate finite bound. No assertion about complexity follows
merely from this decidability argument.
\end{proof}
\end{annexitem}

\begin{annexitem}{annex:p17.htn-summary}{thm:p17-htn-summary}{Exact trace summaries for arbitrary Boolean endpoint and interval constraints}
\begin{proof}
Suppose nonempty traces \(\rho,\eta\) concatenate, so the final state of
\(\rho\) equals the initial state of \(\eta\). The first state and first
post-action state of the concatenation come from \(\rho\); its last
pre-action state and final state come from \(\eta\). Its length-one flag
is false. Directly partitioning the trace states gives
\[
 \begin{aligned}
 A(\rho\eta)&=A(\rho)\cap A(\eta),&
 B(\rho\eta)&=B(\rho)\cap A(\eta),\\
 C(\rho\eta)&=A(\rho)\cap C(\eta),&
 D(\rho\eta)&=B(\rho)\cap C(\eta).
 \end{aligned}
\]
The shared boundary state can occur in both intersections without changing
them. These identities also apply when one or both traces have length one;
the indicated inclusive and exclusive state ranges give the same formulas.
One block needs no concatenation. Repetition gives the summary of any
nonempty list of compatible blocks.

For exact constraint evaluation, name a block's first and last primitive
positions. The total block order determines the order of anchors from
different blocks. Within a block, the first precedes the last unless its
length-one flag is true, when they coincide. Thus first/last expressions and
ordering tests are determined without knowing the full numerical lengths.
The state before or after either anchor is one of the four stored boundary
states.

Consider an interval from an anchor \(e\) to an anchor \(e'\), with the
state-index convention \(\pi(e)\leq j<\pi(e')\). If the endpoints coincide
or occur in reverse order, its universal literal set is all literals. If
they lie in one block in increasing order, they are its first and last
anchors and the required literal set is \(D\) of that block. If they lie
in different blocks in increasing order, intersect the following sets:
\begin{enumerate}
\item from the left block, \(B\) if \(e\) is its first anchor, or the
      literals true in its final state if \(e\) is its last anchor;
\item \(A\) of every intervening block;
\item from the right block, its initial-state literals if \(e'\) is its
      first anchor, or \(C\) if \(e'\) is its last anchor.
\end{enumerate}
These sets cover exactly the states of the half-open interval, allowing
duplicate boundary states. A protection atom holds exactly when its literal
belongs to this intersection. Its negation is evaluated by negating that
membership test, not by replacing the literal with its opposite. All Boolean
combinations are consequently decidable exactly.

An enclosing reference to a compound child is interpreted at its first or
last descendant according to Definition~\ref{def:p17-htn-method}. Repeated
replacement preserves those same extremal anchors. Hence the preceding
evaluation applies both to original first/last expressions and to those
introduced by decomposition; no enclosing interval is discarded.

There are \(2^b\) states, giving at most \(2^{4b}\) boundary-state tuples.
The ground literal set has \(2b\) elements, and each of the four stored
literal sets has at most \(2^{2b}\) possible values. Including the flag
gives the upper bound \(2^{4b}2^{8b}2=2^{12b+1}\). Many such candidates
are unrealizable, but the bound and the composition procedure do not assume
that every candidate is realized.
\end{proof}
\end{annexitem}

\begin{annexitem}{annex:p17.htn-total-order}{thm:p17-htn-total-order}{Finite tabulation with complete boundary and interval information}
\begin{proof}
Ground the finite description. Let \(g\) be the number of ground task
symbols and \(S\leq2^{12b+1}\) the number of candidate trace summaries.
Project states onto the relevant ground atoms. Every precondition, effect
and constraint reads only those coordinates, so projection preserves action
executability, state updates and constraint truth; discarded coordinates
cannot change existence of a solution.
Maintain a table of pairs \((t,\tau)\), initially containing just the
summaries obtained by executing each primitive task from every state in
which it is executable. An initial primitive trace has one action, so its
four boundary states and four literal sets are obtained directly.

For each ground method with head \(t\) and ordered children
\(t_1,\ldots,t_k\), enumerate choices of table entries for the children.
Require compatible consecutive boundary states. Evaluate the method's full
constraint formula by Lemma~\ref{thm:p17-htn-summary}, including its goal
postcondition where applicable. When the formula holds, compose the child
summaries and insert the resulting \((t,\tau)\). Repeat until no new
entry is added. Then treat each ground instance of the initial ordered
network similarly, additionally requiring initial state \(I\). Accept if
any such combination satisfies the initial constraint. Check an empty
initial network directly.

Soundness follows by induction on insertion stages. A primitive entry has
its actual one-action trace. For a newly added entry, choose the witnessing
child traces supplied by the induction hypothesis. Their concatenation is
executable, satisfies the method constraint, and has the computed summary.
The decomposition replacements identify precisely the same endpoints for
enclosing constraints. Therefore the entry witnesses an actual permitted
decomposition with that trace summary.

For completeness, take a finite solution derivation and induct on its
decomposition-tree height. Total ordering of each network forces the
primitive descendants of a child to be a contiguous block before the next
child; descendants of different children cannot interleave. The leaf entries
are initially present. By induction the child summaries are eventually
present, so their ground method instance is considered and its true
constraint is recognized exactly. Its parent entry is inserted. Applying
the same argument to the initial network shows that every solution is
accepted. Arbitrary Boolean and negated interval constraints are covered
by the summary lemma.

There are at most \(gS\) table entries, hence at most \(gS\) productive
rounds. The number of ground methods is at most exponential in a polynomial
in \(L\). Each method has at most \(L\) children, so enumerating summary
tuples contributes at most \(S^L\). Compatibility and formula evaluation
take polynomial time in \(b\) and the grounded description length. The
product of these finite factors is bounded by
\(2^{\operatorname{poly}(L)+O(b\operatorname{poly}(L))}\).
When \(b\) is polynomial this is exponential time; when
\(b\leq2^{\operatorname{poly}(L)}\) it is double exponential time.

The hardness follows from the regular, totally ordered instances constructed
in Annex~\ref{annex:p17.htn-regular}. They are a subclass of the present
instances and already give EXPSPACE-hardness in the lifted case and
PSPACE-hardness in the propositional case. No equality between these lower
and upper bounds is asserted.
\end{proof}
\end{annexitem}

\begin{annexitem}{annex:p17.htn-regular}{thm:p17-htn-regular}{One pending task, bounded-space reachability and machine simulations}
\begin{proof}
\emph{Upper bounds.} Grounding leaves every regular network with at most
one nonprimitive task, ordered relative to its primitive tasks. Guess a
permitted total ordering of those primitive tasks and their ground bindings.
If there is a compound child, its descendants form one block between a
primitive prefix and a primitive suffix: none of the enclosing primitives
can interleave with that block.

Use obligations \((t,\tau)\), with a ground nonprimitive task and a candidate
summary from Definition~\ref{def:p17-htn-trace-summary}. An obligation can
be discharged by a ground method with only primitive children exactly when
their execution has summary \(\tau\) and satisfies the method constraint.
Otherwise a method has one nonprimitive child \(t'\). Guess its summary
\(\tau'\), execute the primitive prefix from the input state of \(\tau\),
require that its resulting state match the input of \(\tau'\), and execute
the primitive suffix from the output of \(\tau'\). Compose the summaries,
check equality with \(\tau\), and evaluate every local constraint exactly
by Lemma~\ref{thm:p17-htn-summary}. A successful check gives a directed edge
\((t,\tau)\to(t',\tau')\). The initial regular network is handled by the
same prefix/child/suffix check from \(I\), or directly if it is primitive.

Every accepted finite path specifies methods ending in a primitive method;
substituting their witnessed traces upwards supplies the required plan.
Conversely, every finite regular derivation has one spine of compound tasks,
so its actual summaries trace such an accepting path. There are at most
\(gS\) obligations. If an accepting path exists, one without a repeated
obligation exists: remove any cycle, retaining the same task and exact
summary at the joining point. This preserves every enclosing constraint.

An obligation, a local method choice and the necessary states/literal sets
use \(O(b+\operatorname{poly}(L))\) space. A path counter bounding search
by \(gS\) uses \(O(\log g+\log S)\) space of the same order. Method
groundings and primitive orderings can be generated as needed. Thus the
problem is in nondeterministic exponential space, or nondeterministic
polynomial space in the propositional case. Savitch's theorem
(Theorem~\ref{thm:p04-savitch}) turns these into EXPSPACE and PSPACE:
squaring an exponential or polynomial space bound stays in the respective
class. With a fixed domain all arities occurring in its schemas are fixed,
so their ground atoms and task instances are polynomially numerous. Any
predicate used only by the initial network is static: after guessing that
network's bindings, its tests can be evaluated from the explicitly given
initial facts and removed. The same polynomial-space procedure therefore
proves the fixed-domain PSPACE upper bound without enumerating all possible
simultaneous groundings of the initial network.

\emph{A fixed-domain PSPACE lower bound.} Fix a deterministic Turing machine
deciding a PSPACE-complete language in polynomial space. Its finite states,
tape alphabet and transition table are fixed. For an input \(x\), use one
constant per cell of its polynomial space bound, with a static successor
relation describing neighboring cells. Unary predicates indicate the head
position and each tape symbol; finitely many nullary predicates indicate
the machine control state. The initial state explicitly describes the
input tape, blanks, initial head and initial control state.
Such a fixed complete-language machine can decide bounded-space universal
acceptance: given a machine description, an input and a unary space bound,
simulate within that bound and stop after the finite configuration bound
if no acceptance occurs. The configuration counter uses polynomial space.
Every polynomial-space language reduces to this problem by inserting its
fixed deciding machine and its polynomial unary bound, proving completeness.

For every machine transition, use a fixed operator schema with variables
for the current and neighboring cells. Its preconditions check the current
control, head, symbol and relevant successor relation. Its effects replace
the scanned symbol, move the head and update the control. A left move uses
the successor relation in the reverse direction; a stay move uses one cell.
The invariant that there is one head, one control state and one symbol per
cell follows immediately from the updates and holds initially. Executable
transition sequences are therefore exactly prefixes of the machine's
computation. The goal is its accepting control state.

The operator set is fixed independently of \(x\); only constants and
initial facts vary, and their size is polynomial in \(\#x\). Apply the
encoding of Theorem~\ref{thm:p17-htn-strips-embedding}. Its methods remain
fixed, regular and totally ordered, use one compound symbol and at most
two tasks per body, and preserve acceptance. Hence a fixed regular domain
is PSPACE-hard. Grounding this bounded-arity family takes polynomial time;
renaming each ground action as a nullary primitive gives propositional
instances with the same restrictions, proving the propositional lower bound.

\emph{An EXPSPACE lower bound with varying arities.} Let a fixed deterministic
machine decide any chosen language in space \(2^{p(\#x)}\), and choose
\(m\geq1\), polynomial in \(\#x\), so that \(2^m\) cells suffice.
Use the two constants 0 and 1. A cell is an \(m\)-tuple of bits, represented
as arguments, not as a function term. Predicates
\(\mathit{Head}(\vec z)\) and \(\mathit{Tape}_a(\vec z)\), for each
tape symbol \(a\), have arity \(m\); control-state predicates are nullary.

The initial description must not explicitly list exponentially many blank
cells. Instead start in an initialization phase with a counter predicate
\(\mathit{Init}(0,\ldots,0)\). For each possible carry position include
an operator whose precondition has the form
\(\mathit{Init}(\vec y,0,1,\ldots,1)\). It writes a blank at that address,
removes that counter fact and adds
\(\mathit{Init}(\vec y,1,0,\ldots,0)\). A final operator at address
\((1,\ldots,1)\) writes its blank and enters the input-loading phase.
These \(m+1\) cases enumerate every bit address exactly once, in binary
order, and no simulation operator is enabled during initialization.

Use one fixed-order loading step per input symbol to replace the blank at
its explicitly given bit address by that symbol. Nullary stage predicates
force these steps in sequence. The last enters the machine's initial control
state and places its head at the first address. Thus all cells have exactly
one symbol, the required initial tape has been produced, and the compact
initialization cannot be bypassed.

For a right-moving machine transition, use one operator per carry position
to change a head address \((\vec y,0,1,\ldots,1)\) to
\((\vec y,1,0,\ldots,0)\), checking the scanned symbol and control and
updating them as above. Left moves use the analogous borrow cases; stay
moves leave the address unchanged. There are no moves beyond the allocated
endpoints. All simulation actions require the appropriate machine control
predicate. Consequently an accepting control state is reachable exactly
when the bounded machine accepts \(x\).

The description uses polynomially many schemas, each with polynomially
many argument positions. The exponentially long initialization is a possible
plan prefix, not part of the encoded input. The operators are function-free
and use positive preconditions; their add/delete effects maintain the stated
machine invariant. Applying the regular HTN encoding again gives one compound
symbol, totally ordered binary method bodies and a one-task initial network.
This is a polynomial reduction for every language in EXPSPACE, proving
EXPSPACE-hardness under all the stated extra restrictions. Combined with
the upper bounds, it completes the three assertions.
\end{proof}
\end{annexitem}

\begin{annexitem}{annex:p17.htn-primitive}{thm:p17-htn-primitive}{Primitive planning: verification and two independent satisfiability reductions}
\begin{proof}
For membership in NP, guess a constant for every variable and a permutation
of the finitely many labelled tasks. Both descriptions have polynomial size.
Simulate the resulting finite action sequence, rejecting an unmet precondition.
Its explicitly stored state contains at most the initial atoms and the atoms
added by these polynomially many operator instances, so this simulation does
not require materializing the exponentially large set of all possible atoms.
Evaluate every equality, ordering, endpoint and interval atom by the stored
trace; interval checks inspect at most the number of task occurrences.
Evaluate the Boolean constraint formula. All these operations take polynomial
time in the input and guessed certificate. With a fixed total order and no
variables, no guess remains, proving the polynomial-time subcase.

For NP-hardness of primitive propositional planning, start with a Boolean
formula \(\beta(p_1,\ldots,p_r)\). Use an initially empty state and two
precondition-free actions per atom: one adds \(p_i\), the other deletes it.
Put every action once in the initial network with no ordering restrictions.
Form the network constraint by replacing each \(p_i\) in \(\beta\) by
the test that \(p_i\) holds after the final occurrence, expressed using
\(\operatorname{last}\) of all labels. A final truth assignment satisfies
the constraint exactly when it satisfies \(\beta\). Conversely, any
satisfying assignment is obtained by ordering each pair with its desired
final effect last; concatenating the pairs realizes all choices together.
This is a polynomial reduction and uses no methods or variables.

For NP-hardness with a totally ordered primitive network, use two constants
\(\mathit{true},\mathit{false}\). For each Boolean variable \(p_i\), use
a distinct unary predicate \(P_i\), one variable \(x_i\), and one
precondition-free action adding \(P_i(x_i)\). Order these actions in a
fixed sequence, starting from the empty state. Replace \(p_i\) in
\(\beta\) by the after-final test \(P_i(\mathit{true})\). Since this
is the only action affecting \(P_i\), that test holds exactly when
\(x_i\) is bound to true. Groundings and Boolean assignments therefore
correspond in both directions. The reduction is polynomial and preserves
total ordering. Unrestricted primitive planning contains both subcases,
so their hardness and the common NP upper bound prove all assertions.
\end{proof}
\end{annexitem}

\begin{annexitem}{annex:p17.htn-strips-embedding}{thm:p17-htn-strips-embedding}{Appending one explicit bookkeeping action}
\begin{proof}
Retain the original constants, predicates, state and action schemas. Introduce
a fresh effect-free primitive action done and one compound task \(t\).
Give \(t\) one terminating method to done. For every original action schema
\(a(\vec x)\), give \(t\) the method consisting of
\(\operatorname{do}[a(\vec x)]\) followed by \(\operatorname{perform}[t]\).
Start with one occurrence \(u:t\); for each goal literal \(\ell\), impose
the after-condition \((u,\ell)\).

Each finite primitive decomposition repeatedly chooses an original action,
then eventually the terminating method. It therefore yields exactly a
sequence \(\sigma\mathit{done}\), with one final done and no other new
actions. The inherited goal condition is tested after this final action.
Because done does not change the state, executable completions are exactly
original executable sequences ending in a goal state, followed by done.
Conversely, any original solution prescribes the successive action methods
and then the terminating method. An empty original solution chooses the
terminating method immediately.

There is one method per original operator and one additional method; their
size is polynomial in the original description. Every method is totally
ordered, contains at most two tasks, and has at most one nonprimitive task,
which is last. Hence the stated regularity and size restrictions hold.
Erasing done gives the promised equality of observable solution sets.
\end{proof}
\end{annexitem}

\begin{annexitem}{annex:p17.htn-expressivity}{thm:p17-htn-expressivity}{Decidability and regular-language obstructions}
\begin{proof}
For any finite function-free STRIPS-style instance, there are finitely many
ground atoms, hence finitely many world states and ground actions. Enumerate
the reachable states from the initial state, applying every executable action
and stopping when no new state appears. A goal state is reachable exactly
when a plan exists. This is a terminating decision procedure, irrespective
of its potentially large resource requirements.

If a computable translation preserved existence of a plan in both directions,
applying this decision procedure after translation would decide general HTN
plan existence. Theorem~\ref{thm:p17-htn-undecidable} excludes that possibility.
The converse direction of existence preservation is necessary: a target
solution not corresponding to any source solution would invalidate this
decidability argument.

For the language claim, make each world state a state of a deterministic
finite automaton, include an additional rejecting dead state, and use ground
actions as alphabet letters. Executable actions follow the planning transition;
undefined actions go to the dead state. The initial automaton state is the
initial world state, and its accepting states are precisely goal states.
The automaton recognizes exactly the STRIPS solution sequences.

Now take two precondition-free, effect-free primitive actions \(a,b\), one
compound task \(t\), and the two totally ordered methods
\(t\) to \(a,b\), and \(t\) to \(a,t,b\). Start with \(t\).
Induction on the number of recursive methods shows that every primitive
derivation yields \(a^m b^m\) for some \(m\geq1\), and choosing the
recursive method \(m-1\) times followed by the base method yields every
such word. Every resulting sequence is executable and all its ordering
constraints hold. Each method has at most one compound child, ordered
relative to every primitive child, so this problem is regular in the source's
definition.

This language is not regular. If a deterministic automaton with \(q\) states
recognized it, two of the \(q+1\) states reached after
\(a,a^2,\ldots,a^{q+1}\) would coincide, at exponents \(i<j\).
Appending \(b^i\) would then either accept both words or reject both.
But \(a^i b^i\) belongs to the language and \(a^j b^i\) does not,
a contradiction.

Renaming letters in a finite automaton, or erasing selected letters by turning
their edges into epsilon edges, yields a finite nondeterministic automaton
for the image language. Such an automaton still recognizes a regular
language. Thus neither symbol renaming nor bookkeeping erasure can make a
STRIPS solution language equal this HTN language. The embedding in the
preceding annex establishes the other direction under its explicit observable
action convention. These are statements about the designated finite,
function-free representations, not every language someone might call planning.
\end{proof}
\end{annexitem}
